\section{The Switch That Turns It Off}
\label{sec:ch09-turning-the-check-off}
\index{satisfiability!disabling the check|(}%

\code{wgc router compose --help} lists six options at 0.129.7, and one of them
takes the subject of section~\ref{sec:ch09-satisfiability} away:

\begin{minted}[fontsize=\footnotesize,breaklines]{text}
--disable-resolvability-validation  This flag will disable the validation for
                                    whether all nodes of the federated graph
                                    are resolvable. Do NOT use unless
                                    troubleshooting.
\end{minted}

The capital letters are theirs. I want to know what a flag does before I decide
whether the warning is proportionate, so I ran it against the unsatisfiable pair
from section~\ref{sec:ch09-satisfiability}: the one that produces four errors
and writes nothing.

With the flag, it composes. Exit code 0, a config on disk, no output but the
success line.

\subsection*{What the config says}

\index{satisfiability!disableEntityResolver@\code{disableEntityResolver}}%
I compared that config with the good one leaf by leaf, and five leaves differ.
Four of the five are the edit echoing back: Mosaic's schema is nineteen
characters longer than it was, so the verbatim copy of it is longer, so its
normalised copy is longer, so the SHA-1 those are stored under is a different
one. Nothing there is about the flag.

The fifth is the only thing anywhere in the file that records the graph is
broken, and it is one property:

\begin{minted}{json}
"keys": [
  { "typeName": "Product", "selectionSet": "id", "disableEntityResolver": true }
]
\end{minted}

That is the honest half. The router is told that Mosaic's \code{Product} key
exists and may not be used to enter it.

The other half is the one that matters, and I checked it in the strongest form I
could rather than eyeballing the two files. The merged client-facing schema in
the broken config and the merged client-facing schema in the good one are the
same string. Not similar: I compared them with \code{===} in node and got
\code{true}.

So the router is handed a schema that advertises \code{price},
\code{availableQuantity}, \code{reviews} and \code{averageRating} on
\code{Product}, and a routing table saying the only service that has them cannot
be reached. Every one of those fields introspects. Every one of them appears in
a client's generated types. Nothing in the schema records that the graph was
composed with the check switched off.

\subsection*{Why that is the worst available outcome}

\index{failures!moved rather than removed}%
A composition failure is the cheapest failure in this book. It happens on a
laptop or in a build, before anything is deployed, in a tool that has both
schemas open in front of it and can name the exact coordinate. Nothing is
running, no client is connected, and the cost of being wrong is that you read a
message and edit a file.

This flag converts that into a runtime failure in a distributed system, and
moves it from the person who caused it to whoever is on call. What the router
actually does when handed this config, I do not know: no router runs in this
chapter, and I am not going to describe behaviour I have not watched. It is
chapter~\ref{ch:enter-the-router}'s first easy experiment, and one request
answers it.

\index{troubleshooting!flags that hide the problem}%
What I will say is that \enquote{do not use unless troubleshooting} is a
strange thing to be able to do at all, and I can construct exactly one case
where I would reach for it: a graph with dozens of subgraphs where the error
names a path I do not believe, and I want to see the routing table the composer
would have produced so I can argue with it. That is a debugging session, and the
config it produces should be deleted at the end of it rather than shipped.

The general shape is worth naming because federation has several switches like
this one. A tool that refuses to produce output is doing the most useful thing
it can do. An override that turns a refusal into a warning is fine; an override
that turns a refusal into silence, and produces an artefact indistinguishable
from a good one, is a way of losing information you cannot get back.

\medskip
Three of the other five options belong in the same paragraph, and none of the
three is exercised here. \code{--suppress-warnings} hides output that Mosaic
never produced, so I have no warning text to show you at all.
\code{--ignore-external-keys} needs \code{@external}, which arrives in
chapter~\ref{ch:entity-resolution-done-right}. \code{--split-configs-enabled}
splits the output across files and wants a router at 0.315.0 or above. The
remaining two are \code{-i} and \code{-o}, which every command in this chapter
has used.

That is composition as a tool. What remains is composition as something a build
does on your behalf, every time, whether or not anybody remembered to ask.

\index{satisfiability!disabling the check|)}%
